\section{Impact and dissemination}
\label{sec:impact}

\subsection{Scientific impact}

If \ref{obj:theory} succeeds, a body of registry epidemiology becomes reinterpretable. Effect
estimates that currently carry an unstated assumption of perfect linkage acquire a stated interval,
and some of them will move. We expect the three reanalyses in \ref{wp:inference} to be the visible
part of that, and the diagnostic package to be the durable part.

If \ref{obj:encoding} succeeds, custodians gain an encoding they can defend rather than one they
hope has not been attacked yet. This matters beyond health data: the same encodings are used in
census linkage, in administrative statistics, and in electoral roll deduplication.

\subsection{Route to use}

\begin{itemize}
  \item All software from \ref{wp:theory}, \ref{wp:encoding}, and \ref{wp:inference} is released
    under a permissive licence at the deliverable date, not at the end of the project.
  \item The protocol specification from \ref{wp:protocol} is published before its implementation,
    so that other groups can implement it independently and so that the security audit has
    something to audit.
  \item The governance model from \ref{wp:governance} is written as a template agreement with
    commentary, because the legal work is the part that other consortia most often cannot afford
    to repeat~\cite{vestergaardnnamdi2023governance}.
\end{itemize}

\subsection{Communication}

Two audiences need different things. Custodians and regulators need the guarantee and the audit
trail, and will be reached through an annual workshop hosted with the three partner custodians.
The public needs to know that a linkage happened, why, and what protects them, and will be reached
through a plain-language annual report and through the lay panel described in Section~\ref{sec:risk},
whose members are recruited from the registry populations rather than from the research community.

\subsection{Beyond the project}

The five-year outcome we are aiming for is narrow and testable: that a registry study submitted in
2033 reports its linkage error the way it currently reports its missing data. That is a change in
convention, and conventions move when the tool is free, the theory is settled, and a custodian will
not release data without it. \proj{} is designed to supply all three.
